\ifdefined\MyWebBook\else
\documentclass[../textbook.tex]{subfiles}
\begin{document}
\fi
\bookchapter{Transformer 网络结构}{ch:transformer}

注意力算子解决位置之间如何交换信息，位置表示则使计算能够区分绝对或相对顺序，但二者尚未构成完整的语言模型。每个位置的表示还需要非线性变换，多层网络需要可训练的梯度路径，输入和输出需要与词元空间连接。Transformer 将这些职责组织为重复的网络层，并通过可见性规则定义信息能够沿什么方向传播。

逐位置前馈网络、残差连接与归一化共同构成 Transformer 层。归一化位置、激活函数和前馈结构决定层内的计算顺序与参数形式；编码器、解码器及编码器—解码器结构则采用不同的信息可见性规则。注意力和位置表示的推导分别见\bookxref{ch:attention}与\bookxref{ch:position-representation}。

\section{Transformer 层的组成}

\subsection{子层分工}
设输入为 $\mat{X}\in\Real^{B\times T\times d}$，其中 $B$ 为批大小，$T$ 为序列长度，$d$ 为隐藏维数。\term{多头注意力}{Multi-Head Attention}{MHA}在自注意力模式下将不同位置的信息汇入当前表示；逐位置前馈网络只沿特征轴计算，所有位置共享同一组参数。两者交替堆叠，使跨位置交互得到的特征能够被非线性重组，再影响下一层的读取。

残差流保持形状 $[B,T,d]$，成为各子层共同读写的表示空间。注意力的输出投影和前馈网络的下投影都必须回到维数 $d$，才能与输入逐元素相加。残差连接采用逐元素相加，通道拼接和跨层特征平均都不属于这一操作。

\begin{table}[htbp]
\centering
\caption{Transformer 核心数学符号与张量形状}\label{tab:transformer-notations}
\begin{tabularx}{\linewidth}{p{24mm}X}
\toprule 符号 & 含义与形状\\\midrule
$B,\mat{S},T$ & 批大小、源序列长度、目标序列长度。\\
$d,h,d_h$ & 隐藏维数、注意力头数、单头维数；标准等宽结构满足 $d=hd_h$。\\
$f$ & 前馈网络中间维数，不必固定为 $4d$。\\
$N_e,N_d$ & 编码器层数、解码器层数。\\
$\mat{X},\mat{C}$ & 目标状态 $[B,T,d]$、源端记忆 $[B,\mat{S},d]$。\\
$\mat{V}_{\mathrm{vocab}}$ & 输出词表大小，与注意力值矩阵 $\mat{V}$ 区分。\\
$\vect{\gamma},\vect{\beta}$ & 归一化增益与偏置，均为 $d$ 维向量。\\
$\mat{J}_F$ & 将张量展平后，函数 $F$ 相对于输入的雅可比。\\
\bottomrule
\end{tabularx}
\end{table}

\begin{assumption}[结构推导的共同假设]
所有子层在有效位置保持隐藏维数；标准多头结构满足 $d=hd_h$；注意力每个有效查询至少有一个可见键；LayerNorm 与 RMSNorm 沿特征轴计算，除显式手算例外均使用 $\epsilon>0$。掩码与输入位置均视为给定的非训练状态。
\end{assumption}

矩阵公式使用行向量约定，即 $\mat{X}\mat{W}$；求导时为方便讨论，将单个位置视为列向量，雅可比按“输出分量对输入分量”排列。二者只是记号选择，参数存储布局必须在实现中对应。除明确说明外，结构公式省略 Dropout，其影响在训练部分单独讨论。

\subsection{注意力子层的接口}
令查询侧输入为 $\mat{X}_q\in\Real^{T\times d}$，键值侧输入为 $\mat{X}_k\in\Real^{S\times d}$。第 $r$ 个头计算
\begin{align}
 \mat{Q}_r&=\mat{X}_q\mat{W}_{Q,r},&\mat{K}_r&=\mat{X}_k\mat{W}_{K,r},&\mat{V}_r&=\mat{X}_k\mat{W}_{V,r},\\
 \mat{A}_r&=\softmax(\frac{\mat{Q}_r\mat{K}_r^{\mathrm{T}}}{\sqrt{d_h}}+\mat{M}),&&
 \mat{O}_r=\mat{A}_r\mat{V}_r.
\end{align}
其中三个投影矩阵为 $d\times d_h$，$\mat{M}$ 为 $T\times \mat{S}$ 加性掩码，Softmax 沿键轴归一化。拼接各头再乘 $\mat{W}_O\in\Real^{d\times d}$，得到 $T\times d$ 输出。自注意力令两侧输入相同；交叉注意力则允许两侧来自不同网络、具有不同长度。

“自注意力中 $\mat{Q}=\mat{K}=\mat{V}$”通常只是对接口输入来源的简写，投影后的三个矩阵一般互不相等。三个投影可以打包为一次矩阵乘法，参数区间仍然相互独立，没有发生权重共享。交叉注意力的查询与键值来自不同张量，一个输入无法无条件完成全部三组投影。

\section{逐位置前馈网络}

\subsection{两层前馈网络的非线性作用}
\term{逐位置前馈网络}{Feed-Forward Network}{FFN}定义为
\begin{equation}\label{eq:tr-ffn}
 F(\mat{X})=\phi(\mat{X}\mat{W}_1+\vect{b}_1)\mat{W}_2+\vect{b}_2,
 \quad \mat{W}_1\in\Real^{d\times f},\quad \mat{W}_2\in\Real^{f\times d}.
\end{equation}
$\vect{b}_1\in\Real^f,\vect{b}_2\in\Real^d$ 按位置广播，张量路径为 $[B,T,d]\to[B,T,f]\to[B,T,d]$。隐藏宽度 $f$ 控制可用中间特征数量；它是结构超参数，不由注意力头数决定。

如果去掉激活函数，两个仿射变换可以合成一个仿射变换：
\begin{equation}
 (\mat{X}\mat{W}_1+\vect{b}_1)\mat{W}_2+\vect{b}_2=\mat{X}(\mat{W}_1\mat{W}_2)+(\vect{b}_1\mat{W}_2+\vect{b}_2).
\end{equation}
因此仅增加中间宽度并不能提供同样的非线性表达能力。激活使不同输入采用不同的有效特征组合。例如 ReLU 把负预激活置零，其局部雅可比相当于在两个投影之间插入与输入有关的对角选择矩阵。

把无偏置 FFN 写为逐中间特征的和，可得到
\begin{equation}
 F(\vect{x})=\sum_{j=1}^{f}\phi(\vect{x}\vect{w}_{1,j})\vect{w}_{2,j}^{\mathrm{T}},
\end{equation}
$w_{1,j}$ 为上投影第 $j$ 列，$w_{2,j}^{\mathrm{T}}$ 为下投影第 $j$ 行。上投影矩阵提取输入特征的基底投影，下投影矩阵对激活后的特征进行重组，激活函数则根据输入动态控制各基底通道的增益。神经元在深层表征中常呈现出多义性（Polysemanticity）与特征叠加（Superposition），若要分析单个神经元的概念对应关系，需借助稀疏自编码器（SAE）等特征解离工具进行实证分析。

\subsection{ReLU 与 GELU}
原始 Transformer 使用\term{整流线性单元}{Rectified Linear Unit}{ReLU} $\operatorname{ReLU}(u)=\max(0,u)$\cite{vaswani2017attention}。其正半轴导数为一，负半轴导数为零，零点不可微，实现需采用约定的次梯度。ReLU 的输出非负，但经过下投影与残差相加后，最终表示仍可具有任意符号。

\term{高斯误差线性单元}{Gaussian Error Linear Unit}{GELU}使用标准正态分布函数 $\Phi$\cite{hendrycks2016gelu}：
\begin{equation}
 \begin{aligned}
  \operatorname{GELU}(u)&=u\Phi(u),\\
  \Phi(u)&=\int_{-\infty}^{u}\frac{\conste{}^{-\frac{t^2}{2}}}{\sqrt{2\constpi{}}}\dif t.
 \end{aligned}
\end{equation}
其导数由乘积规则得到
\begin{equation}
 \operatorname{GELU}'(u)=\Phi(u)+u\frac{\conste{}^{-\frac{u^2}{2}}}{\sqrt{2\constpi{}}}.
\end{equation}
与 ReLU 在负区间直接截断不同，GELU 对负输入仍可能输出较小负值。例如 $u=-1$ 时输出约为 $-0.1587$，$u=1$ 时约为 $0.8413$，零点导数为 $0.5$。

一种常用近似为
\begin{equation}
 \operatorname{GELU}(u)\approx\frac u2
 \left[1+\tanh\left(\sqrt{\frac2\constpi{}}(u+0.044715u^3)\right)\right].
\end{equation}
精确形式与近似形式不逐位等价。迁移检查点或核对数值时，需要匹配所采用的定义，配置中出现的“GELU”这一名称本身不足以确定数值。

\subsection{GLU、GEGLU 与 SwiGLU}
\term{门控线性单元}{Gated Linear Unit}{GLU}将两条投影分支逐元素相乘。省略偏置时，可以统一写为
\begin{equation}\label{eq:tr-glu}
 F_{\mathrm{gate}}(\mat{X})=[g(\mat{X}\mat{W}_g)\odot(\mat{X}\mat{W}_u)]\mat{W}_d,
 \quad \mat{W}_g,\mat{W}_u\in\Real^{d\times f},\quad \mat{W}_d\in\Real^{f\times d}.
\end{equation}
经典 GLU 的门函数为 Sigmoid；\term{高斯误差线性门控单元}{GELU-Gated Linear Unit}{GEGLU}使用 GELU；\term{Swish门控线性单元}{Swish-Gated Linear Unit}{SwiGLU}采用 Swish 门，本文取其参数为一，即\term{Sigmoid线性单元}{Sigmoid Linear Unit}{SiLU}，写为 $\operatorname{SiLU}(u)=u\sigma(u)$，其中 $\sigma(u)=\frac{1}{1+\conste{}^{-u}}$\cite{shazeer2020glu}。门控分支与内容分支都由当前输入决定，最后仍通过下投影回到残差宽度。

\begin{figure}[htbp]
\centering
\includegraphics[width=.9\linewidth]{\subfix{../../figures/theory/revision-ch08-gated-ffn.pdf}}
\caption[标准 FFN 与门控 FFN 结构对比]{标准 FFN 与门控 FFN 结构对比。门控结构新增输入投影，并通过逐元素乘法形成输入相关的交互。}\label{fig:tr-ffn}
\end{figure}

图\ref{fig:tr-ffn}中的两条分支在同一中间维数上对齐。对单个中间分量，若 $a=xw_g,\vect{b}=xw_u$，输出乘积为 $g(a)\vect{b}$，其微分为
\begin{equation}
 \dif[g(a)\vect{b}]=\vect{b}\,g'(a)\dif a+g(a)\dif \vect{b}.
\end{equation}
两条分支的梯度分别依赖另一条分支的输出。Sigmoid 门取值在零与一之间；SiLU 门可以为负且没有有限正上界，因此 SwiGLU 的门值表示乘法调制系数，不能解释为开启概率。稀疏专家选择还需要独立的路由机制。

无偏置两层 FFN 的参数量为 $2df$，同宽门控 FFN 为 $3df$。公平比较时应区分保持中间维数与保持参数预算：若原 FFN 宽度为 $4d$，等参数门控宽度约为 $\frac{8d}{3}$，因为 $3d(\frac{8d}{3})=2d(4d)$。实际维数还可能为硬件对齐取整。门控 FFN 的线性投影数量已经改变，仅更换式\eqref{eq:tr-ffn}中的激活名称无法得到它。表\ref{tab:mlp-variants}对比了不同前馈网络变体的参数量与计算特性。

\begin{table}[htbp]
\centering
\caption{前馈网络变体结构与参数量对比}\label{tab:mlp-variants}
\begin{tabularx}{\linewidth}{p{20mm}p{28mm}p{18mm}p{22mm}X}
\toprule
前馈网络类型 & 门控激活函数 & 投影矩阵数 & 等参数隐藏宽度 & 机制特性与代表制品 \\\midrule
标准两层 FFN & 无门控（ReLU / GELU） & 2 ($\mat{W}_1, \mat{W}_2$) & $f \approx 4d$ & 经典两层升维降维变换，参数量 $8d^2$（BERT、GPT-2/3） \\
GEGLU & $\operatorname{GELU}(\mat{X}\mat{W}_g) \odot \mat{X}\mat{W}_u$ & 3 ($\mat{W}_g, \mat{W}_u, \mat{W}_d$) & $f \approx \frac{8}{3}d$ & 双路并行投影并经 GELU 门控调制，门控分支调节特征幅度（T5、PaLM） \\
SwiGLU & $\operatorname{SiLU}(\mat{X}\mat{W}_g) \odot \mat{X}\mat{W}_u$ & 3 ($\mat{W}_g, \mat{W}_u, \mat{W}_d$) & $f \approx \frac{8}{3}d$ & 采用平滑、非单调的自门控激活（LLaMA-2/3、Qwen2.5/3、DeepSeek-V3） \\
\bottomrule
\end{tabularx}
\end{table}

\section{残差连接}

\subsection{残差连接的梯度传播}
将任一保持形状的子层写为 $F$，残差更新为
\begin{equation}
 \vect{x}_{\ell+1}=\vect{x}_\ell+F_\ell(\vect{x}_\ell).
\end{equation}
当 $F_\ell$ 接近零时，该层接近恒等映射。网络不必让一个全新的复合变换从随机初始化中学习如何复制所有输入信息，而可以围绕已有表示学习修正。

对输入求导得到
\begin{equation}
 \begin{aligned}
  \mat{J}_\ell&=\mat{I}+\mat{J}_{F_\ell},\\
  \frac{\partial\mathcal L}{\partial \vect{x}_\ell}
  &=(\mat{I}+\mat{J}_{F_\ell})^{\mathrm{T}}
  \frac{\partial\mathcal L}{\partial \vect{x}_{\ell+1}}.
 \end{aligned}
\end{equation}
上游梯度既直接传回，也经过子层雅可比传回。堆叠多层后，端到端雅可比是这些矩阵的有序乘积；把它改写成标量乘积会丢失各层的输入依赖。

\subsection{恒等路径的稳定性边界}
即使每层都有 $\mat{I}$ 项，梯度仍可能放大或衰减。一维构造中，若 $F_\ell(\vect{x})=a\vect{x}$，则 $L$ 层导数为 $(1+a)^L$；$a=0.1$ 时随深度增长，$a=-0.1$ 时随深度衰减，$a=-1$ 时甚至直接为零。残差提供可用路径，参数、尺度、归一化和优化过程共同决定实际稳定性。

残差流方差随输入与增量的相关性变化。若输入与增量近似不相关，则相加后的方差约为两者方差之和；若相关，还包含协方差项。因此深层网络常配合归一化、初始化或残差缩放。采用 $\vect{x}_{\ell+1}=\vect{x}_\ell+\alpha_\ell F_\ell(\vect{x}_\ell)$ 会把雅可比改为 $\mat{I}+\alpha_\ell \mat{J}_{F_\ell}$，缩放影响前向和反向，属于对网络函数的实质修改。

\section{归一化}

\subsection{LayerNorm 的统计轴}
\term{层归一化}{Layer Normalization}{LayerNorm}按一个位置的特征维计算统计量\cite{ba2016layernorm}。对 $\vect{x}\in\Real^d$，定义
\begin{align}
 \mu&=\frac1d\sum_{i=1}^{d}x_i,&
 \vect{c}&=\vect{x}-\mu\mathbf1,&
 s&=\sqrt{\frac1d \vect{c}^{\mathrm{T}}\vect{c}+\epsilon},\\
 \widehat{\vect{x}}&=\frac{\vect{c}}s,&
 \operatorname{LN}(\vect{x})&=\vect{\gamma}\odot\widehat{\vect{x}}+\vect{\beta}.&&
\end{align}
$\epsilon>0$ 防止零方差造成除零，$\vect{\gamma},\vect{\beta}\in\Real^d$ 是可训练向量。这里除以 $d$，与用于估计总体方差的无偏修正 $d-1$ 不同。归一化的目标是变换当前向量，统计估计则另有其口径。

LayerNorm 不在批轴或时间轴聚合，因此批大小和其他样本不会直接改变当前位置的统计量。但此前注意力若错误读取填充位置，污染已经进入当前向量，LayerNorm 并不能消除这种信息泄漏。它依赖当前样本的即时统计量，与使用训练阶段移动统计量的典型 BatchNorm 也不同。

\subsection{不变性与边界}
对整体平移 $\vect{x}+a\mathbf1$，中心化结果不变，所以 LayerNorm 精确消除该方向。忽略 $\epsilon$ 且方差非零时，正比例缩放也被消除；有非零 $\epsilon$ 时只在其相对方差足够小时近似成立。负比例缩放会改变中心化向量的符号，输出与正比例缩放不同。

以 $\vect{x}=(1,3)$ 为例，在 $\vect{\gamma}=(1,1),\vect{\beta}=(0,0)$ 下，均值为二，方差为一，输出为 $\frac{-1,1}{\sqrt{1+\epsilon}}$。$\vect{x}'=(11,13)$ 的输出相同。加上可学习仿射变换后，输出未必仍为零均值或单位方差；$\vect{\gamma}$ 的各分量不同，$\vect{\beta}$ 也可以引入偏移。

\subsection{LayerNorm 雅可比推导}
令中心化矩阵 $\mat{P}=\mat{I}-\frac{\mathbf1\mathbf1^{\mathrm{T}}}{d}$，则 $\vect{c}=\mat{P}\vect{x}$、$\mat{P}^{\mathrm{T}}=\mat{P}$、$\mat{P}\vect{c}=\vect{c}$。有
\begin{equation}
 \begin{aligned}
  \dif \vect{c}&=\mat{P}\dif \vect{x},\\
  \dif s&=\frac{\vect{c}^{\mathrm{T}}\dif \vect{x}}{ds}.
 \end{aligned}
\end{equation}
对 $\widehat{\vect{x}}=\frac{\vect{c}}{s}$ 应用商法则，得到
\begin{equation}
 \begin{aligned}
  \dif\widehat{\vect{x}}&=
  \left(\frac{\mat{P}}s-\frac{\vect{c}\vect{c}^{\mathrm{T}}}{ds^3}\right)\dif \vect{x},\\
  \mat{J}_{\mathrm{LN}}&=\operatorname{diag}(\vect{\gamma})
  \left(\frac{\mat{P}}s-\frac{\vect{c}\vect{c}^{\mathrm{T}}}{ds^3}\right).
 \end{aligned}
\end{equation}
均值和尺度计算使同一位置的各特征相互耦合，雅可比因而不是简单的对角缩放。设输出上游梯度为 $\vect{g}$，令 $\vect{h}=\vect{\gamma}\odot \vect{g}$，则输入梯度也可写为
\begin{equation}\label{eq:tr-lnback}
 \nabla_{\vect{x}}\mathcal L=\frac1s
 \left[\vect{h}-\overline h\mathbf1-\widehat{\vect{x}}\,
 \overline{\vect{h}\odot\widehat{\vect{x}}}\right],
 \quad \overline h=\frac1d\sum_i h_i.
\end{equation}
仿射参数梯度分别为 $\vect{g}\odot\widehat{\vect{x}}$ 和 $\vect{g}$，训练时再对批及位置求和。式\eqref{eq:tr-lnback}说明某一特征上的损失会通过统计量影响其他特征，均值与方差因而是参与求导的中间量。

\subsection{RMSNorm 的尺度归一化}
\term{均方根归一化}{Root Mean Square Normalization}{RMSNorm}保留均值，仅用均方根缩放\cite{zhang2019rmsnorm}：
\begin{equation}
 \begin{aligned}
  r&=\sqrt{\frac1d \vect{x}^{\mathrm{T}}\vect{x}+\epsilon},\\
  \operatorname{RMSNorm}(\vect{x})&=\vect{\gamma}\odot\frac{\vect{x}}r.
 \end{aligned}
\end{equation}
这里采用无偏置版本，参数量为 $d$；LayerNorm 的增益与偏置共为 $2d$。RMSNorm 不具有整体平移不变性。对 $(1,3)$，其输出为 $\frac{1,3}{\sqrt{5+\epsilon}}$，与 LayerNorm 的 $\frac{-1,1}{\sqrt{1+\epsilon}}$ 明显不同。

相同的微分方法给出
\begin{equation}
 \mat{J}_{\mathrm{RMS}}=\operatorname{diag}(\vect{\gamma})
 \left(\frac{\mat{I}}r-\frac{\vect{x}\vect{x}^{\mathrm{T}}}{dr^3}\right).
\end{equation}
LayerNorm 多出中心化投影 $\mat{P}$，RMSNorm 则保留均值方向。二者采用不同的统计量，前向与梯度性质随之不同。归一化的计算成本只是整网开销的一部分，其吞吐收益还取决于其他算子与访存。表\ref{tab:norm-comparison}对比了主流特征归一化方案的数学形式与系统特性。

\begin{table}[htbp]
\centering
\caption{主流特征归一化方案对比}\label{tab:norm-comparison}
\begin{tabularx}{\linewidth}{p{20mm}p{32mm}p{16mm}p{16mm}X}
\toprule
归一化方案 & 归一化公式 & 均值中心化 & 参数量 & 优化特性与工程优势 \\\midrule
LayerNorm & $\vect{\gamma}\odot\frac{\vect{x}-\mu}{\sqrt{\sigma^2+\epsilon}}+\vect{\beta}$ & 是 (减 $\mu$) & $2d$ & 经典平移和缩放不变性，需两遍访存并保存均值中间量 \\
RMSNorm & $\vect{\gamma}\odot\frac{\vect{x}}{\sqrt{\frac{1}{d}\|\vect{x}\|^2+\epsilon}}$ & 否 (保留均值) & $d$ & 消除均值计算，减少访存与显存开销，训练稳定性与 LayerNorm 相当，为现代大模型首选 \\
DeepNorm & 残差缩放型 LayerNorm & 是 & $2d$ & 理论保障千层极深网络无梯度爆炸，多见于超大规模多语言模型 \\
\bottomrule
\end{tabularx}
\end{table}

\section{Pre-LN 与 Post-LN 的结构差异}

\subsection{子层形式与归一化位置}
对一个注意力或 FFN 子层 $F$，\term{后置层归一化}{Post-Layer Normalization}{Post-LN}将归一化放在残差相加之后；\term{前置层归一化}{Pre-Layer Normalization}{Pre-LN}将其放在子层输入端：
\begin{align}
 \vect{y}_{\mathrm{post}}&=\operatorname{LN}(\vect{x}+F(\vect{x})),\\
 \vect{y}_{\mathrm{pre}}&=\vect{x}+F(\operatorname{LN}(\vect{x})).
\end{align}
同样的 $F$ 和同样的归一化参数通常产生不同结果。Post-LN 的残差分支最终也经过归一化，Pre-LN 的主残差流则绕过本子层归一化。归一化与残差相加的顺序因此是网络层定义的一部分。

\begin{figure}[htbp]
\centering
\includegraphics[width=.9\linewidth]{\subfix{../../figures/theory/revision-ch08-norm-paths.pdf}}
\caption[归一化位置对前向与反向路径的影响]{归一化位置对前向与反向路径的影响。Pre-LN 的恒等项绕过本子层归一化，末端归一化仍影响最终损失。}\label{fig:tr-norm}
\end{figure}

沿图\ref{fig:tr-norm}的路径应用链式法则，二者雅可比分别为
\begin{align}
 \mat{J}_{\mathrm{post}}&=\mat{J}_{\mathrm{LN}}(\vect{x}+F(\vect{x}))[\mat{I}+\mat{J}_F(\vect{x})],\\
 \mat{J}_{\mathrm{pre}}&=\mat{I}+\mat{J}_F(\operatorname{LN}(\vect{x}))\mat{J}_{\mathrm{LN}}(\vect{x}).
\end{align}
Pre-LN 的主残差流具备\emph{显式恒等路径}，浅层表征与深层梯度可直接贯通；Post-LN 的两条局部反向路径均须\emph{连续乘以归一化雅可比矩阵}，导致反向梯度范数随网络深度加深产生扰动与退化\cite{xiong2020layernorm}。因此 Post-LN 通常严格依赖学习率预热（Warmup）与精细初始化以维持初期优化稳定。

\begin{warning}[归一化位置与权重兼容性]
Pre-LN 与 Post-LN 具备完全不同的前向计算图与梯度传播动力学。从预训练模型转换架构时，无法直接复用既有权重；即便张量形状严格匹配，改变归一化相对位置也会导致输出表征彻底失真。
\end{warning}

\subsection{末端归一化}
典型 Pre-LN 堆叠在最后添加一次归一化，再进入输出头，以控制主残差流到预测层的尺度。层内归一化作用于子层输入，末端归一化作用于整个堆叠的输出，两者的位置和作用对象不同。它的雅可比也参与从损失到主干的梯度传播。

本章将原论文配置记为“Post-LN、LayerNorm、ReLU”，将一个便于比较的变体记为“Pre-LN、LayerNorm、GELU、末端 LayerNorm”。RMSNorm 与 SwiGLU 是另外两项独立设计选择，不由 Pre-LN 自动推出。逐项改变配置才能判断差异来源；一次同时替换三项设置后观察到损失变化，其来源需要进一步拆分。

\section{Transformer 的架构类型}

\subsection{编码器与因果解码器}
\term{仅编码器}{Encoder-only}{}结构堆叠自注意力和 FFN，通常允许有效位置双向读取，输出上下文表示，之后连接分类、标注或检索任务头。\term{仅解码器}{Decoder-only}{}结构通常堆叠因果自注意力和 FFN，再预测下一词元，通常已经删除交叉注意力子层，与原始解码器的组成有出入。

网络架构的拓扑形态由子层模块组成与注意力可见性掩码共同决定：\emph{仅编码器}面向双向表征抽取，\emph{仅解码器}依托因果掩码成为现代生成式大语言模型的主流标准，\emph{编码器—解码器}则依托交叉注意力保留了显式的源到目标记忆交互通道。因此，不能仅凭名称指代其计算路径，需明确其掩码拓扑与子层配置。

\subsection{编码器—解码器的数据流}
\term{编码器—解码器}{Encoder--Decoder}{}先将源输入 $x_{1:S}$ 编码为记忆 $\mat{C}\in\Real^{B\times S\times d}$；解码器根据右移后的目标前缀查询这份记忆，输出 $[B,T,d]$ 表示。交叉注意力的查询来自目标侧，键和值来自源侧，权重形状为 $[B,h,T,\mat{S}]$。

\begin{figure}[!tbp]
\centering
\includegraphics[width=.94\linewidth]{\subfix{../../figures/theory/revision-ch08-architecture.pdf}}
\caption[原始 Transformer 编解码架构]{原始 Transformer 编解码架构。主数据流自下而上；灰色旁路显式表示残差连接，加法后执行 LayerNorm。蓝色跨列箭头将最终源记忆送入每个解码器层。依据原始结构绘制\cite{vaswani2017attention}。}\label{fig:tr-architecture}
\end{figure}

阅读图\ref{fig:tr-architecture}时，先沿源侧观察记忆的形成，再沿目标侧观察三个子层如何更新目标表示。交叉注意力的查询来自目标侧，因此输出长度为 $T$；键和值来自源侧，其长度为 $\mat{S}$。训练时，右移操作使位置 $t$ 的输入不包含待预测的 $y_t$，因果掩码再阻止该位置读取后续目标输入。图中各子层均先残差相加、后归一化。

图中的上下游连接需要由层内公式具体化。以原始 Post-LN 顺序定义编码器第 $\ell$ 层，且各归一化参数独立：
\begin{align}
 \mat{U}_\ell&=\operatorname{LN}_{\ell,1}
 [\mat{C}_{\ell-1}+\operatorname{MHA}_{e,\ell}(\mat{C}_{\ell-1},\mat{C}_{\ell-1})],\\
 \mat{C}_\ell&=\operatorname{LN}_{\ell,2}[\mat{U}_\ell+F_{e,\ell}(\mat{U}_\ell)].
\end{align}
这里 MHA 的两个输入分别表示查询侧与键值侧。源嵌入加位置表示得到 $\mat{C}_0$，最终 $\mat{C}=\mat{C}_{N_e}$。不同层通常具有独立参数；“堆叠相同结构”说的是结构形式，各层权重并不复用。

对应解码器层包含三个子层：
\begin{align}
 \mat{U}_\ell&=\operatorname{LN}_{\ell,1}
 [\mat{X}_{\ell-1}+\operatorname{MHA}_{c,\ell}(\mat{X}_{\ell-1},\mat{X}_{\ell-1})],\\
 \mat{Z}_\ell&=\operatorname{LN}_{\ell,2}
 [\mat{U}_\ell+\operatorname{MHA}_{x,\ell}(\mat{U}_\ell,\mat{C})],\\
 \mat{X}_\ell&=\operatorname{LN}_{\ell,3}[\mat{Z}_\ell+F_{d,\ell}(\mat{Z}_\ell)].
\end{align}
下标 $c$ 表示目标侧因果自注意力，$x$ 表示交叉注意力。每个目标层使用同一份最终源记忆，但各层的交叉注意力投影参数独立。源序列不随目标生成增长，因此编码器可以在生成前运行一次。

Pre-LN 解码器则可写为
\begin{align}
 \mat{U}_\ell&=\mat{X}_{\ell-1}+\operatorname{MHA}_{c,\ell}
 (\operatorname{LN}_{\ell,1}(\mat{X}_{\ell-1}),\operatorname{LN}_{\ell,1}(\mat{X}_{\ell-1})),\\
 \mat{Z}_\ell&=\mat{U}_\ell+\operatorname{MHA}_{x,\ell}(\operatorname{LN}_{\ell,2}(\mat{U}_\ell),\mat{C}),\\
 \mat{X}_\ell&=\mat{Z}_\ell+F_{d,\ell}(\operatorname{LN}_{\ell,3}(\mat{Z}_\ell)).
\end{align}
这里明确选择只归一化交叉注意力查询侧输入，$\mat{C}$ 已由源端网络输出；若另一配置还归一化源记忆，应另作定义，实现中额外增加的运算需要在定义里写明。

\subsection{编码器、因果与前缀掩码对照}
令 $A_{ij}=1$ 表示位置 $i$ 允许读取位置 $j$，并暂时不含填充。对长度四的序列，双向、因果和前两位置为条件前缀的掩码分别为
\begin{equation}
 \mat{A}_{\mathrm{bi}}=\begin{bmatrix}1&1&1&1\\1&1&1&1\\1&1&1&1\\1&1&1&1\end{bmatrix},\quad
 \mat{A}_{\mathrm{causal}}=\begin{bmatrix}1&0&0&0\\1&1&0&0\\1&1&1&0\\1&1&1&1\end{bmatrix},\quad
 \mat{A}_{\mathrm{prefix}}=\begin{bmatrix}1&1&0&0\\1&1&0&0\\1&1&1&0\\1&1&1&1\end{bmatrix}.
\end{equation}
\term{前缀解码器}{Prefix Decoder}{}的条件部分可双向交互，生成部分读取全部前缀及生成历史。设前缀长度为 $P$，其规则为
\begin{equation}
 A_{ij}=\mathbf1[(i\le P\land j\le P)\lor(i>P\land j\le i)].
\end{equation}
训练若把前缀内部各位置也按普通下一词元目标评分，可能发生泄漏，因为前缀状态可读取前缀未来；应按实际目标选择监督位置。前缀双向条件化与独立编码器加交叉注意力是两种结构，信息变换与参数形式都不同。

第四种情况是交叉注意力，其矩阵为 $T\times \mat{S}$，一般允许每个目标查询读取全部有效源位置。它不在源轴应用目标侧因果三角形。填充、文档隔离或流式源输入可以进一步限制各类掩码，但应在任务定义中明确。

\subsection{多层因果性定理}
如果初始位置 $i$ 的状态只依赖输入 $x_{\le i}$，且每层位置 $i$ 只读取前层 $j\le i$ 的状态，那么下一层位置 $i$ 也只依赖 $x_{\le i}$。逐位置 FFN、残差和特征轴归一化都不会引入其他时间位置，因此归纳可知任意层均满足因果性。

编码器—解码器中，同样的结论是目标位置仅依赖源输入及真实目标前缀，未来目标不进入该位置的计算。这个证明的前提包括正确右移：若把当前位置真实目标直接作为同位置输入，即使使用下三角掩码，也已在初始表示中泄漏答案。这里的因果性专指生成顺序上的信息约束。

\section{模型规模计算}

\subsection{单层形状表}
标准等宽多头注意力要求 $d=hd_h$。现代主流大语言模型普遍采用图\ref{fig:transformer-decoder-block}所示的纯解码器（Decoder-Only）前置归一化结构。

\begin{figure}[htbp]
\centering
\includegraphics[width=0.92\linewidth]{\subfix{../../figures/theory/transformer-decoder-block.pdf}}
\caption[现代纯解码器架构与张量流向]{现代纯解码器架构与张量流向。采用 Pre-RMSNorm 残差连接主干，注意量子层整合 GQA 与 RoPE 旋转位置编码，前馈子层采用 SwiGLU 门控结构。}\label{fig:transformer-decoder-block}
\end{figure}

以下以自回归注意力与普通前馈网络为例，偏置按末轴广播，各算子张量路径如表\ref{tab:transformer-tensor-shapes}所示。

\begin{table}[htbp]
\centering
\caption{单层 Transformer 算子张量形状}\label{tab:transformer-tensor-shapes}
\begin{tabularx}{\linewidth}{p{28mm}p{44mm}X}
\toprule 运算 & 输入或参数 & 输出\\\midrule
QKV 投影 & $[B,T,d]$，三个 $d\times d$ 矩阵 & 三个 $[B,T,d]$\\
拆头 & $d=hd_h$ & 三个 $[B,h,T,d_h]$\\
缩放点积 & $\frac{\mat{Q}\mat{K}^{\mathrm{T}}}{\sqrt{d_h}}$ & $[B,h,T,T]$\\
归一化与读取 & 掩码后 Softmax，再乘 $\mat{V}$ & $[B,h,T,d_h]$\\
拼头与输出投影 & 拼接末轴，乘 $d\times d$ & $[B,T,d]$\\
残差与归一化 & 两条 $[B,T,d]$ 支路 & $[B,T,d]$\\
FFN 上投影与激活 & $d\times f$ & $[B,T,f]$\\
FFN 下投影 & $f\times d$ & $[B,T,d]$\\
词表输出头 & $d\times \mat{V}_{\mathrm{vocab}}$ & $[B,T,\mat{V}_{\mathrm{vocab}}]$\\
\bottomrule
\end{tabularx}
\end{table}
交叉注意力将键值长度换为 $\mat{S}$，分数形状为 $[B,h,T,\mat{S}]$，输出仍为目标侧 $[B,T,d]$。其输出长度由查询数决定，不由源长度决定。

\subsection{参数量核算}
带偏置的标准 MHA 具有四个 $d\times d$ 投影及四个 $d$ 维偏置，共 $4d^2+4d$ 个参数；带偏置普通 FFN 为 $2df+f+d$；两个 LayerNorm 共为 $4d$。因此，编码器层或仅解码器层的参数量为
\begin{equation}
 P_{\mathrm{layer}}=4d^2+2df+f+9d.
\end{equation}
含交叉注意力的解码器层有两组 MHA、一个 FFN 和三个 LayerNorm，总计
\begin{equation}
 P_{\mathrm{encdec\ layer}}=8d^2+2df+f+15d.
\end{equation}
两式不含嵌入、位置参数、末端归一化及输出头，也不适用于直接删偏置、共享键值头或门控 FFN 的配置。

例如 $d=512,f=2048$ 时，前者为 $\num{3152384}$，后者为 $\num{4204032}$。六个编码器层与六个解码器层合计 $\num{44138496}$ 个参数，但这不是模型总参数量。若使用无偏置 SwiGLU 和两次无偏置 RMSNorm，仅解码器层对应为 $4d^2+3df+2d$，比较时必须同时说明投影是否含偏置。

\subsection{权重共享的前提与梯度}
若目标嵌入为 $\mat{E}\in\Real^{V_{\mathrm{vocab}}\times d}$，输出头可以共享其转置：
\begin{equation}
 \vect{z}_t=\vect{h}_t\mat{E}^{\mathrm{T}}+\vect{b}_{\mathrm{out}}.
\end{equation}
独立输出矩阵被省去，减少 $\mat{V}_{\mathrm{vocab}}d$ 个参数。共享矩阵收到输入查表和输出投影两条路径的梯度之和。输出路径的 Softmax 通常使许多词表行得到梯度，因而某行在嵌入查表中被设为不更新，它在共享输出路径中仍可能发生变化。

源嵌入、目标嵌入和输出头是否能够进一步共享，需要兼容的词表编号及维数。仅仅词表大小相同不够，同一行必须指向同一词元。共享降低参数自由度，输入表示与输出语义是否一致仍取决于训练目标和数据。

\subsection{资源开销}
标准自注意力投影规模为 $\mat{O}(BTd^2)$，两次注意力矩阵乘法为 $\mat{O}(BT^2d)$，普通 FFN 为 $\mat{O}(BTdf)$。显式保存逐头权重需要 $\mat{O}(BhT^2)$ 元素。交叉注意力包含 $\mat{O}(BT\mat{S}d)$ 的读取计算，并有源侧和目标侧投影。

保持总宽度 $d$ 不变而增加头数，标准投影总参数量并不会线性增加；每头维数相应缩小。然而显式注意力矩阵含头轴，存储开销可能改变。高效内核可以避免完整物化权重矩阵，这改变的是访存和执行路径；在数学定义上，密集注意力仍然包含全部位置对之间的交互。

\begin{example}
以下构造一个两位置、单头、隐藏维数为二的因果 Pre-LN 层，使用普通 ReLU FFN。输入已包含嵌入和位置表示：
\begin{equation*}
 \mat{X}=\begin{bmatrix}1&3\\4&2\end{bmatrix},\qquad
 \mat{W}_Q=\mat{W}_K=\mat{W}_V=\mat{W}_O=\mat{I}_2.
\end{equation*}
所有偏置为零，所有归一化增益为一，省略 Dropout。本例各次方差非零，为获得简洁精确值取 $\epsilon=0$；实际实现应保持正 $\epsilon$ 并相应计算数值。

前馈矩阵取
\begin{align*}
 \mat{W}_1&=\begin{bmatrix}1&0&1\\0&1&-1\end{bmatrix},\\
 \mat{W}_2&=\begin{bmatrix}1&0\\0&1\\1&-1\end{bmatrix}.
\end{align*}
层后再做一次最终 LayerNorm，并采用三词元输出头
\begin{equation*}
 \mat{W}_{\mathrm{out}}=\begin{bmatrix}1&0&-1\\0&1&1\end{bmatrix},
\end{equation*}
三列依次对应词元 $\mat{A},B,\mat{C}$。本例使用独立输出矩阵，不进行权重共享。

逐行归一化得到
\begin{equation*}
 \widehat{\mat{X}}=\operatorname{LN}(\mat{X})=\begin{bmatrix}-1&1\\1&-1\end{bmatrix}.
\end{equation*}
因投影为恒等矩阵，$\mat{Q}=\mat{K}=\mat{V}=\widehat{\mat{X}}$。缩放分数与因果权重为
\begin{align*}
 \mat{S}&=\frac{\mat{Q}\mat{K}^{\mathrm{T}}}{\sqrt2}
 =\begin{bmatrix}\sqrt2&-\sqrt2\\-\sqrt2&\sqrt2\end{bmatrix},\\
 \mat{A}&=\begin{bmatrix}1&0\\1-p&p\end{bmatrix},
 \quad p=\frac{\conste{}^{\sqrt2}}{\conste{}^{\sqrt2}+\conste{}^{-\sqrt2}}.
\end{align*}
$p\approx0.944193$。第一行只允许一个键，故权重严格为一；第二行比较两个键。记 $c=2p-1\approx0.888386$，则
\begin{align*}
 \mat{O}&=\mat{A}\mat{V}=\begin{bmatrix}-1&1\\c&-c\end{bmatrix},\\
 \mat{U}&=\mat{X}+\mat{O}=\begin{bmatrix}0&4\\4+c&2-c\end{bmatrix}.
\end{align*}
注意力改变了表示，输出维数仍为二，因而可以与输入相加。

$\mat{U}$ 两行的特征大小关系没有逆转，因此
\begin{equation*}
 \operatorname{LN}(\mat{U})=\begin{bmatrix}-1&1\\1&-1\end{bmatrix}.
\end{equation*}
进入 FFN 后，
\begin{align*}
 \operatorname{LN}(\mat{U})\mat{W}_1&=\begin{bmatrix}-1&1&-2\\1&-1&2\end{bmatrix},\\
 \operatorname{ReLU}(\operatorname{LN}(\mat{U})\mat{W}_1)&=\begin{bmatrix}0&1&0\\1&0&2\end{bmatrix},\\
 F(\operatorname{LN}(\mat{U}))&=\begin{bmatrix}0&1\\3&-2\end{bmatrix}.
\end{align*}
第二条残差得到
\begin{equation*}
 \mat{H}=\mat{U}+F(\operatorname{LN}(\mat{U}))
 =\begin{bmatrix}0&5\\7+c&-c\end{bmatrix}.
\end{equation*}
此处 FFN 的中间维数为三，输出重新回到二。若遗漏下投影，则无法进行这次残差相加。

末端归一化后，$\overline{\mat{H}}=[(-1,1);(1,-1)]$。于是词表分数为
\begin{equation*}
 \mat{Z}=\overline{\mat{H}} \mat{W}_{\mathrm{out}}=
 \begin{bmatrix}-1&1&2\\1&-1&-2\end{bmatrix}.
\end{equation*}
逐行 Softmax 得到近似概率
\begin{equation*}
 \mat{P}=\begin{bmatrix}
 0.035119&0.259496&0.705385\\
 0.843795&0.114195&0.042010
 \end{bmatrix}.
\end{equation*}
若两个监督目标分别为 $B,\mat{A}$，平均负对数似然为
\begin{equation*}
 \mathcal L=-\frac12[\log(0.259496)+\log(0.843795)]\approx0.759429.
\end{equation*}
至此，输入、归一化、注意力、两条残差、FFN、末端归一化、词表投影和损失已经形成完整计算链。

二维且 $\epsilon=0$ 的 LayerNorm 把任意两个不相等分量归一化为相同或相反的符号模式，这使本例中若干归一化结果重复。这一性质由二维和零 $\epsilon$ 的设定共同产生；增加维数、使用正 $\epsilon$ 或改变归一化位置，都会改变数值路径。
\end{example}

\section{模型执行接口}

\subsection{分类概率模型}
Transformer 特征主干接收连续表示并返回连续表示。完整语言模型还需要分词器、嵌入、位置机制、输出词表层、目标函数与生成循环。输出头将 $[B,T,d]$ 隐藏表示映射为词表分数，再经归一化得到词元概率。

训练时，源文本 $x$ 和真实目标 $y$ 定义条件似然
\begin{equation}
 p_\theta(y\mid x)=\prod_{t=1}^{\mathrm{T}}p_\theta(y_t\mid y_{<t},x).
\end{equation}
\term{教师强制}{Teacher Forcing}{}使用真实目标前缀，右移后的目标输入与因果掩码使各位置可并行计算。目标为 $(\mat{A},B,\mathrm{EOS})$ 时，输入为 $(\mathrm{BOS},\mat{A},B)$，同位置输出分别监督三个目标。填充标签应排除，提示或条件前缀是否监督则由目标定义决定。

生成时用模型已选词元替代真实前缀。编码器记忆可复用，解码器逐步追加目标词元，并在 EOS 或长度上限处终止。缓存避免重复计算历史键值，自回归决策依赖仍然逐步串行。缓存正确性要求位置编号、掩码、精度和参数与完整前向一致。

\begin{algorithm}[教师强制的训练前向]
\AlgoInput{源词元、包含 EOS 的目标词元、模型参数和有效目标长度。}
\AlgoOutput{有效目标位置上的平均损失。}
\AlgoState{源记忆、目标层表示和受监督位置集合。}
\begin{myalgorithmic}[1]
\State 构建源序列填充掩码，计算嵌入与位置表示，并编码为源端记忆
\State 将目标序列右移一位构建解码器输入：起始词元置于首位，其余位置放入前一真实目标词元；保持位置级输入--目标配对
\State 构建目标端因果掩码、目标填充掩码与交叉注意力源端掩码，随后按序执行全部解码层
\State 执行指定的最终归一化与词表投影；仅在选定目标位置累加交叉熵，并按有效位置数归一化
\State 处理完全部目标位置后结束；若无有效监督位置，按指定规则跳过批次或报告状态，不在零分母上求平均
\end{myalgorithmic}
\textbf{不变量：}位置 $t$ 的预测只依赖源输入与 $y_{<t}$；显式右移与库内部的下一词元配对只能采用一套，重复移位会改变监督对象。
\end{algorithm}

\begin{algorithm}[条件生成的自回归执行]
\AlgoInput{源序列、启动与终止词元、最大生成长度、固定模型和词元选择策略。}
\AlgoOutput{每条样本的生成词元序列。}
\AlgoState{源记忆、已生成前缀、可选键值缓存、样本终止标志。}
\begin{enumerate}
\item 关闭训练随机性，编码源序列一次；将目标前缀初始化为启动词元，各样本标为未终止。
\item 对未终止样本计算当前最后有效目标位置的分布。使用缓存时同步更新位置编号、可见性和当前层键值。
\item 按规定策略选择下一词元并追加到前缀；生成 EOS 的样本标为终止，后续不再追加普通生成内容。
\item 若全部样本终止或达到长度上限，返回结果；否则回到步骤二。
\end{enumerate}
\textbf{不变量：}缓存内容与当前模型参数及前缀对应；终止样本不再改变其已生成答案。模型前向、词元选择和停止策略分别承担不同职责。
\end{algorithm}

\subsection{正则化与评价口径}
训练中的\term{随机失活}{Dropout}{}可以作用于注意力权重、FFN 中间表示或残差增量。按保留概率校正的 Dropout 在期望上保持对应分量的尺度，但一次采样会改变前向输出；注意力权重经过 Dropout 后，单次实现的行和也不必仍为一。进行确定性结构对照时必须关闭随机性，或对齐全部随机状态。

标签平滑将独热目标替换为 $q=(1-\eta)e_y+\eta u$，$u$ 为规定的参考分布。损失为 $-\sum_vq_v\log p_v$，与普通真实标签负对数似然不同。若训练记录平滑损失而验证记录未平滑损失，二者差值还包含平滑本身带来的偏移。

训练集用于参数更新，验证集用于选择检查点和停止时机，测试集用于最终评价。精确匹配会把合法同义表达判错，字符或词片段指标也不能覆盖全部语义质量；评价指标与划分应在专门评估中明确。注意力热力图呈现读取模式，因果归因还需干预或对照实验。

\subsection{模型结构的兼容性}
从原理实现迁移至工程框架（如 PyTorch、Megatron 或 vLLM）时，需依次核验参数权重映射、投影矩阵分块布局、归一化位置、激活函数实现、偏置项使能、掩码极性定义以及末端输出头投影。对于同构模型，应在固定随机种子与相同输入张量下，在前向 Logits 与反向梯度上分别核对绝对误差与相对误差。

布尔掩码的真值方向是接口约定，其含义随具体接口而定。将“允许读取”矩阵转换为“禁止读取”矩阵时需要取反；加性掩码则直接作用于分数。优化提示也代替不了对实际可见性矩阵的核验。错误提示可能使后端选择不适用的执行路径，故接口封装应显式保存掩码语义。

模型制品还包含非训练参数状态，例如固定位置编码、频率、最大长度及归一化常数。保存方式可以是这些状态本身，也可以是按固定配置精确重建的规则；重载一致性要求可训练权重与这些状态共同匹配。现成翻译模型的分词器、特殊词元、位置机制与权重是一组配套对象，与另一个原理模型的整数输入不能混用。

因果性检验可固定参数与随机状态后扰动未来词元，验证早期输出严格保持不变；填充隔离检验需改变 PAD 对应内容并核对有效序列输出不变；数值归一化检验则应覆盖恒等常数输入与趋近零方差边界。单元级一致性检验保障算子实现正确，模型任务能力则需通过评测基准与服务基准负载综合测量。






\chapterreferences
\myendchapterdocument
